\section{From solving an equation to repeating a map}\label{ch12:sec:iteration}
Section~\ref{ch04:sec:unique} proved that the equation
\[
x=\frac{x+1}{2}
\]
has exactly one real solution, namely $x=1$. Section~\ref{ch05:sec:recurrence} reached the same number in a different way. Starting from $x_0=0$, it applied the rule $x\mapsto(x+1)/2$ again and again and obtained
\[
0,\ \tfrac12,\ \tfrac34,\ \tfrac78,\ \ldots,\qquad\text{in general }x_n=1-2^{-n},
\]
and Section~\ref{ch11:sec:convergence} proved that these numbers converge to $1$. So the number $1$ can be seen in two ways: as the solution of an equation, and as the limit of a repeated update. This chapter turns the second point of view into a general method for solving equations.

The general setting is a map $T:X\to X$, where $X$ is a metric space as in Chapter~\ref{ch:limits}. (Recall from Section~\ref{ch03:sec:functions} that a map is the same thing as a function.) A \emph{fixed point} of $T$ is a point $x_*\in X$ with $T(x_*)=x_*$, that is, a point that the map leaves where it is. Section~\ref{ch05:sec:recurrence} already used this name for the number $1$ and the rule $x\mapsto(x+1)/2$. Starting from any point $x_0\in X$, we can apply $T$ over and over:
\[
x_1=T(x_0),\qquad x_2=T(x_1),\qquad x_3=T(x_2),\qquad\ldots,
\]
in general $x_{n+1}=T(x_n)$ for every $n\in\NN$. The points $x_0,x_1,x_2,\ldots$ are called the \emph{iterates} of $T$ starting from $x_0$. If the iterates converge and their limit is a fixed point, then repeating the map is a way of solving the equation $T(x)=x$.

Repeating a map is a very simple algorithm, but it does not always work. Here are four ways in which it can fail.
\begin{itemize}
\item The iterates can run away. For $T(x)=2x$ on $\RR$, starting from $x_0=1$, the iterates are $1,2,4,8,\ldots$, although $T$ has the fixed point $0$.
\item The iterates can jump back and forth forever. For $T(x)=1-x$ on $\RR$, starting from $x_0=0$, the iterates are $0,1,0,1,\ldots$. The fixed point is $1/2$, since $1-\tfrac12=\tfrac12$, but every iterate stays at distance $1/2$ from it.
\item There can be several fixed points, so that ``the solution'' is not well defined. For $T(x)=x^2$ on the interval $[0,1]$, both $0$ and $1$ are fixed points.
\item The iterates can approach a point that is missing from the space. For $T(x)=x/2$ on the open interval $(0,1)$, starting from $x_0=1/2$, the iterates are $\tfrac12,\tfrac14,\tfrac18,\ldots$. They approach $0$, which does not belong to $(0,1)$, and in fact $T$ has no fixed point in $(0,1)$ at all.
\end{itemize}
We want conditions on $T$ and $X$ that rule out all of these failures in advance, so that we do not have to compute a few iterates and hope that the pattern continues. The first condition says that $T$ brings points closer together.

\par\noindent\begin{minipage}{\linewidth}
\begin{definition}[Contraction]\label{ch12:def:contraction}
Let $(X,d)$ be a metric space, and let $q$ be a real number with $0\le q<1$. A map $T:X\to X$ is a \emph{contraction with factor $q$} if
\[
d(Tx,Ty)\le q\,d(x,y)\qquad\text{for all }x,y\in X.
\]
This inequality is called the \emph{contraction inequality}, and $q$ is called a \emph{contraction factor} of $T$.
\end{definition}
\end{minipage}\par

To save brackets we write $Tx$ for $T(x)$; both mean the same point. In words, the definition says that after applying $T$, any two points are at most the fraction $q$ of their original distance apart. With $q=1/2$, for instance, every pair of points ends up at most half as far apart as before. If $T$ is a contraction with factor $q$, it is also one with every larger factor below one, which is why we say ``a'' contraction factor rather than ``the'' contraction factor.

\subsection*{Reading the contraction inequality}
Two features of the definition matter in everything that follows.

First, $x$ and $y$ are \emph{any} two points of $X$. They need not be consecutive iterates, and neither of them needs to be close to a fixed point. For $T(x)=(x+1)/2$ on $\RR$, with the usual distance $d(x,y)=|x-y|$, we compute
\[
|T(x)-T(y)|=\left|\frac{x+1}{2}-\frac{y+1}{2}\right|
=\left|\frac{(x+1)-(y+1)}{2}\right|
=\frac12|x-y|.
\]
So this $T$ is a contraction with factor $1/2$. The constant terms cancel when two outputs are compared, just as they did in the uniqueness proof of Section~\ref{ch04:sec:unique}.

Second, one number $q$ must work for all pairs at once. It is not enough that every pair of different points moves closer together. If the amount of shrinking were allowed to depend on the pair, the ratio $d(Tx,Ty)/d(x,y)$ could creep up towards $1$ as the points change, and then the conclusions of this chapter may fail. Section~\ref{ch12:sec:hypotheses} gives a map that brings every pair of different points strictly closer together and still has no fixed point.

\subsection*{The outputs must be allowed inputs}
The arrow in $T:X\to X$ is part of the definition: every output $T(x)$ must again be a point of $X$. A map from a set to itself is called a \emph{self-map}. This condition is what allows us to apply $T$ to $x_1$, then to $x_2$, and so on, exactly as Section~\ref{ch03:sec:functions} required the outputs of one function to be allowed inputs of the next before composing them.

The condition is easy to overlook when a formula makes sense for more numbers than we intend to use. Take $T(x)=(x+1)/2$ again, but now only on the interval $[0,1/2]$. The contraction inequality still holds for points of this interval, with factor $1/2$, because it holds for all real numbers. Yet $T(1/2)=3/4$ lies outside $[0,1/2]$, so $T$ is not a map from $[0,1/2]$ to itself. Starting from $x_0=0$, we get $x_1=1/2$ and then $x_2=3/4$, and the iterates have already left the interval. On $[0,1]$, by contrast, $T$ is a self-map: if $0\le x\le1$, then $1\le x+1\le2$, so $1/2\le T(x)\le1$.

So before using a contraction, we check two separate things: that $T$ maps $X$ into $X$, and that the contraction inequality holds for all pairs of points of $X$.

\pauseq{Is $T(x)=x^2$ a contraction of the interval $[0,1]$? Is it a contraction of the interval $[0,1/3]$?}{On $[0,1]$ it is not. It is a self-map there, because $0\le x\le1$ gives $0\le x^2\le1$. But for $x=1$ and $y=1/2$ we get $|T(x)-T(y)|=1-\tfrac14=\tfrac34$, which is larger than $|x-y|=\tfrac12$; no factor $q<1$, and not even $q=1$, can work for this pair. This agrees with the list above, where the same map had two fixed points. On $[0,1/3]$ it is a contraction. It is a self-map, because $0\le x\le1/3$ gives $0\le x^2\le1/9\le1/3$. For the contraction inequality, factor the difference of two squares: $|x^2-y^2|=|x-y|\,|x+y|$. For $x,y\in[0,1/3]$ we have $0\le x+y\le2/3$, so $|T(x)-T(y)|\le\tfrac23|x-y|$, and $q=2/3$ is a contraction factor.}
